% Capítulo 1 — Régua conceitual-normativa

\section{\MakeUppercase{Transição energética e seus critérios de ruptura}}

\subsection{Racionalidade ambiental como critério epistêmico}

Conforme \citeonline{leff2006}, a racionalidade ambiental constitui-se como contracorrente da racionalidade econômica dominante e opõe, à redução da natureza a capital fungível, a reapropriação social do mundo natural por meio dos saberes locais.

A incomensurabilidade é o mecanismo central que \citeauthoronline{leff2006} identifica para distinguir as duas racionalidades. Pois constata-se que os valores simbólicos, ecológicos, epistemológicos e políticos da natureza e da cultura resistem à homologação ao equivalente universal do mercado.

\subsection{Decrescimento e bem viver}

Sob outro ângulo, \citeauthoronline{latouche2009} sustenta que o decrescimento não constitui antimodernismo, mas recusa do imaginário econômico segundo o qual o crescimento se apresenta como única forma social aceitável.

Em chave complementar, \citeauthoronline{acosta2016} opera no plano relacional, propondo o bem viver como suficiência relacional fixada pela comunidade em sua relação com o território.
